% !TEX root = ../Thesis.tex

%---------------------------------------------------------------------------%
%->> Frontmatter
%---------------------------------------------------------------------------%
%-
%-> 生成封面
%-

\maketitle% 生成中文封面
\MAKETITLE% 生成英文封面
%-
%-> 作者声明
%-
\makedeclaration% 生成声明页
%-
%-> 中文摘要
%-
\intobmk\chapter*{摘\quad 要}% 显示在书签但不显示在目录
\setcounter{page}{1}% 开始页码
\pagenumbering{Roman}% 页码符号

多模光纤因其能够支持数千个空间模式,在紧凑且微创的形态下实现高分辨率成像,成为生物医学内窥成像领域的理想探针。然而,光在多模光纤中传输时,模间耦合、模式干涉以及侧壁反射会导致复杂的散斑图案,严重限制了图像传输的保真度。近年来,深度学习技术为多模光纤成像提供了新的解决方案,通过端到端的监督学习建立散斑图案与目标图像之间的映射关系,实现了高质量的图像重建。然而,如何进一步提升多模光纤的信息容量和重建保真度,充分利用光纤支持的数千个模式,仍是亟待解决的关键问题。

本文系统研究了基于深度学习的多模光纤散斑图像重建方法,重点探索了通过多维度编码策略提升信息传输容量。首先基于公开的多芯光纤散斑数据集验证了U-Net神经网络的有效性,在手写数字和服装数据集上分别实现了0.982和0.922的平均结构相似性指数(SSIM)。在此基础上,提出了同时利用全息编码和偏振编码的双编码方案:采用空间光调制器加载二值振幅全息图设计25个全息标签,结合三种偏振态编码,理论上可实现75路复用传输。针对双编码散斑重建任务,提出了改进的DeepLeakyU-Net网络架构,在编码器路径采用Leaky ReLU激活函数并使用步长卷积替代池化层,在三种偏振态下均实现了平均SSIM达0.93的高保真度重建,并通过无监督聚类分析验证了不同编码条件下散斑图案的可区分性。

本文的主要创新点包括:(1)首次将全息编码与偏振编码相结合,提出双编码方案,显著提升多模光纤的信息传输容量;(2)提出改进的DeepLeakyU-Net网络架构,提升了散斑图像的重建性能;(3)采用远场成像机制,利用阶跃型多模光纤的角度相关性实现基于角域信息的图像重建;(4)通过无监督聚类分析验证了编码方案的鲁棒性和可区分性。

本研究为多模光纤成像系统的性能提升提供了新的技术路径,在生物医学内窥成像等领域具有重要的应用前景。未来工作可进一步优化编码参数,探索更高复用密度,并将方法应用于动态环境下的实时成像场景。

\keywords{多模光纤,全息编码,偏振编码,深度神经网络,图像重建}% 中文关键词
%-
%-> 英文摘要
%-
\intobmk\chapter*{Abstract}% 显示在书签但不显示在目录

Multimode fibers (MMFs) have emerged as ideal probes for biomedical endoscopic imaging due to their capability to support thousands of spatial modes while maintaining a compact and minimally invasive form factor for high-resolution imaging. However, optical transmission through MMFs is complicated by intermodal coupling, modal interference, and sidewall reflections, which collectively generate complex speckle patterns and severely restrict image transmission fidelity. In recent years, deep learning techniques have provided novel solutions for MMF imaging by establishing end-to-end mappings between speckle patterns and target images through supervised learning, achieving high-quality image reconstruction. Nevertheless, how to further enhance the information capacity and reconstruction fidelity of MMFs while fully exploiting the thousands of modes supported by the fiber remains a critical challenge.

This dissertation systematically investigates deep learning-based speckle image reconstruction methods for multimode fibers, focusing on enhancing information transmission capacity through multi-dimensional encoding strategies. The effectiveness of U-Net neural networks was first validated on publicly available multi-core fiber speckle datasets, achieving average Structural Similarity Index Measure (SSIM) of 0.982 and 0.922 for handwritten digit and fashion datasets respectively. Building upon this, a dual-encoding scheme was proposed that simultaneously exploits holographic and polarization encoding: a spatial light modulator displays binary amplitude holograms with twenty-five holographic labels arranged in a $5\times5$ grid, combined with three polarization states, theoretically enabling 75 multiplexed transmission channels. For dual-encoded speckle reconstruction, an improved DeepLeakyU-Net architecture was proposed, adopting Leaky ReLU activation functions in the encoder path and strided convolutions instead of pooling layers, achieving high-fidelity reconstruction with an average SSIM of 0.93 across all three polarization states, with unsupervised clustering analysis validating the discriminability of speckle patterns under different encoding conditions.

The main innovations of this work include: (1) the first combination of holographic and polarization encoding through a dual-encoding scheme, significantly enhancing the information transmission capacity of MMFs; (2) the proposal of an improved DeepLeakyU-Net architecture that enhances speckle image reconstruction performance; (3) the adoption of a far-field imaging mechanism that exploits angular correlations in step-index MMFs to achieve angular-domain-based image reconstruction; (4) validation of the robustness and discriminability of the encoding scheme through unsupervised clustering analysis.

This research provides a novel technical approach for enhancing the performance of multimode fiber imaging systems, with significant application prospects in biomedical endoscopic imaging and related fields. Future work may further optimize encoding parameters, explore higher multiplexing densities, and apply the methods to real-time imaging scenarios under dynamic conditions.

\KEYWORDS{Multimode fiber, Holographic encoding, Polarization encoding, Deep neural network, Image reconstruction}% 英文关键词

\pagestyle{enfrontmatterstyle}%
\cleardoublepage\pagestyle{frontmatterstyle}%

%---------------------------------------------------------------------------%
